% !TEX root = ../main.tex
\chapter{How to detect gravitational waves}\label{ch:c28}
\begin{paracol}{2}

\begin{sources}
\begin{tabularx}{\linewidth}{@{}l l L@{}}
\toprule
\textbf{Source} & \textbf{Length} & \textbf{Content}\\
\midrule
\seg{28}{1}{Lecture 28: How to detect gravitational waves} & 1:21:43 & Lecturer: Bernard F.\ Schutz, not Schuller. The principle of detection; bar detectors and the Weber controversy; beam detectors and the three term relation; spacecraft tracking, pulsar timing, one arm detectors with clocks, atom interferometry; why two arms; the short arm limit; the cosmic microwave background (BICEP2); laser interferometers: noise, pendulum suspensions, stored light, power and signal recycling, two misconceptions corrected; quantum noise and squeezed light; LISA Pathfinder and the transponder. Student questions and the closing of the school.\\
\bottomrule
\end{tabularx}

\smallskip
\textbf{Transcript} \file{materials/transcripts/C28-V01.txt} (YouTube captions, lower case, many mis-hearings). The lecture was given in 2015, before the first detection. It continues chapter~\ref{ch:c27}; it is the last lecture of the school.
\end{sources}

\begin{coursenote}[Lecturer]
Lecturer: Bernard F.\ Schutz, not Schuller. The formula of the three term relation was on a slide and is only described in the captions: it is given below in the structure the lecturer describes, with the coefficients marked [?].
\end{coursenote}

\begin{reading}[Where this chapter sits in the other notes]
No other Schuller set has a matching chapter. In this book: sources and amplitudes are chapter~\ref{ch:c27}; linearised gravitational waves and the transverse traceless gauge are chapters~\ref{ch:c24} and~\ref{ch:c25} (not written yet); geodesics, proper time and the equivalence principle are chapters~\ref{ch:c07} and~\ref{ch:c13}.
\end{reading}

\section{The questions of principle}\label{sec:c28-intro}

The lecturer wants to address the issues of principle people always ask: ``if gravitational waves are a distortion of space, and all of space is distorted, and the light is distorted, how can you detect gravitational waves with light?'' The aim is a clear physical understanding of what happens in detection; then the different ways of detecting gravitational waves (not only interferometers), some technologies and, if time allows, the space technologies. ``This is not primarily about technology, it is about the whole enterprise of detecting gravitational waves and how it fits into physics'' \ts{28}{1}{0:00}\ts{28}{1}{0:21}\ts{28}{1}{0:41}\ts{28}{1}{1:01}.

A reminder of the polarization diagram of the last lecture, drawn differently: the two orthogonal polarizations with the stretching along the $x,y$ axes and rotated by $45^\circ$, and the tidal representation. If general relativity were not correct there could be other polarizations, but general relativity is correct to about $1\,\%$ in its predictions for gravitational waves, from the binary pulsar observations: a good basis to describe detection; small changes at the $1\,\%$ level will teach things but not affect whether the waves are detected \ts{28}{1}{1:23}\ts{28}{1}{1:43}\ts{28}{1}{2:04}. The interferometer again: with the arrangement of a dark port, no light without a gravitational wave, light when there is one \ts{28}{1}{2:25}.

\section{Bar detectors and the Weber controversy}\label{sec:c28-bars}

The first attempt to detect gravitational waves was with \textbf{bar detectors}: a picture of Joseph Weber, the pioneer of the field, with his first bar at the University of Maryland (probably around 1965). Weber considered interferometry but the laser technology was not there (the first continuous wave laser: 1961 ``maybe'', as the lecturer asks the audience), he rejected it and did what he could \ts{28}{1}{2:45}\ts{28}{1}{3:05}\ts{28}{1}{3:26}.

\paragraph{How a bar works.} A metal cylinder (aluminium or steel), operating mainly at its \emph{resonant frequency}. A gravitational wave comes in from the side; the stretching in the $x$ direction stretches the bar, which then oscillates after the wave has passed; strain gauges at the point of maximum strain, the centre of the bar, are supposed to tell about the excitation of the fundamental normal mode. For a steel or aluminium bar about a metre long the frequency is around a kilohertz: the most sensitive frequency \ts{28}{1}{3:47}\ts{28}{1}{4:08}\ts{28}{1}{4:28}.

\paragraph{The inherent limit.} To take advantage of the tidal effect one wants bigger (``bigger is better''), but it becomes impractical: the longest bar made for this purpose is about $3$ m, and the larger the bar the \emph{lower} the resonant frequency, being an acoustic oscillation: a kilometre long piece of aluminium would have a far too low frequency, so one cannot take advantage of the tidal effect. Bars were pushed to about $10^{-19}$ by cooling to very close to absolute zero (liquid helium temperature), ``still not good enough for any detection'': bars are not operated any more; there have been a few projects for more imaginative bars, but ``it is an inherently limited technology'' \ts{28}{1}{4:49}\ts{28}{1}{5:10}\ts{28}{1}{5:30}\ts{28}{1}{5:51}.

\begin{history}[The Weber controversy]
As said in class (not checked): in the late 1960s Weber believed he was seeing gravitational waves and published it: he had two bars, one in Maryland and one in Chicago, believed they were excited in coincidence, and showed that as the Earth turned more coincidences occurred when the bar was broadside onto the centre of the galaxy, so the excitations seemed to come from there, ``even more convincing''. But his sensitivity was low and nobody could understand how there could be so many events of such high amplitude from the centre of the galaxy (the galaxy would lose a lot of mass over millions of years). His data analysis was suspect: not digital, looking with the eye at two chart recordings, ``you can fool yourself, find things that you want to find even if you tell yourself you are being honest''. The only way to test it was to build other bar detectors, and most of the groups now in the interferometry business started out building bars to test Weber; all came up negative; Weber continued to see waves, tried to invent new kinds of physics to explain it, was on the fringe and ``died eventually a rather broken, very sad individual''. The consequence is tremendous caution: LIGO sets a very high standard of confidence for the first detection, ``we don't want to fool ourselves, we don't want people in the community asking us whether we are fooling ourselves: Weber has left a legacy that still affects the field today''.
\end{history}
\ts{28}{1}{5:51}\ts{28}{1}{6:12}\ts{28}{1}{6:33}\ts{28}{1}{6:53}\ts{28}{1}{7:14}\ts{28}{1}{7:34}\ts{28}{1}{7:54}\ts{28}{1}{8:15}\ts{28}{1}{8:36}

\section{Beam detectors: the interaction of light with a gravitational wave}\label{sec:c28-beam}

The most effective detectors are \textbf{beam detectors}: they use the interaction of a beam of light with the gravitational wave field \ts{28}{1}{8:56}.

In a local inertial frame (freely falling, with a Minkowski metric), the gravitational wave is a small perturbation on the Minkowski metric, a linearized wave; it travels at the speed of light through the frame, ``always locally at the speed of light''; but globally gravitational waves distort the geometry of spacetime and this affects the \emph{travel time} between one place and another. Beam detectors are based on measuring the travel time of light through a gravitational wave field \ts{28}{1}{9:17}\ts{28}{1}{9:38}\ts{28}{1}{9:58}.

\paragraph{The thought experiment.} A beam of light is sent from a freely falling source (``really in outer space with nothing around, that is what we will do with LISA; not what we do on the ground, because on the ground we are not freely falling, so one has to ask about the next order'') to a mirror, and reflects directly back. How much \emph{proper time} elapses on a good special-relativity clock carried by the source for the return trip? If gravitational waves go through, the return time changes: ``that is the principle of measurement used in interferometers and other beam detectors'' \ts{28}{1}{10:18}\ts{28}{1}{10:39}\ts{28}{1}{10:59}.

This is \emph{independent of coordinate systems and gauges}: a physical object sends out light, another physical object reflects it, both in free fall, and a clock measures proper time locally: ``nothing to do with coordinates; you can calculate in any coordinate system in general relativity, and the return time will be the same, a coordinate independent measurable'' \ts{28}{1}{11:19}\ts{28}{1}{11:40}.

\begin{theorem}[The three term relation, structure as described in the lecture]\label{thm:c28-three}
Let a plane gravitational wave with polarization amplitude $h_+$ cross the light path of length $L$, with $\theta$ the angle between the wave fronts and the path of the light (the lecturer: ``the angle between the direction of travel of the gravitational wave and that of the light''; the two descriptions differ by $\pi/2$, see the coursenote below). The rate of change of the return time $T$ (proper time on the local clock) is
\[
  \frac{\mathrm dT}{\mathrm dt}\;=\;1+c_1(\theta)\,h_+(t)+c_2(\theta)\,h_+\bigl(t+L(1-\sin\theta)\bigr)+c_3(\theta)\,h_+(t+2L)\qquad[c_i\ \text{not on the captions: [?]}]
\]
(the $1$ is the unperturbed return time: for $h_+=0$ there is no variation, and the interest is in the small variations of order $h$). Here $t$ is the time at which the light is \emph{sent}. The remarkable thing: only three times of the waveform matter, \emph{the emission time $t$, the intermediate time $t+L(1-\sin\theta)$ at the mirror, and the return time $t+2L$}; ``it is not a continuous time, it is just three different points in the evolution of that waveform that affect the return time''. The return time is itself changed by $h$, but that is already a second order effect inside the arguments.
\end{theorem}
\ts{28}{1}{12:01}\ts{28}{1}{12:21}\ts{28}{1}{12:42}\ts{28}{1}{13:03}\ts{28}{1}{13:44}\ts{28}{1}{14:05}\ts{28}{1}{14:25}\ts{28}{1}{14:46}\ts{28}{1}{15:07}\ts{28}{1}{15:27}\ts{28}{1}{15:48}\ts{28}{1}{16:08}

(Called the ``three term relation'', ``a very prosaic name for the fundamental formula for this whole business''. The lecturer takes the $x$ axis in the plane common to the wave and the light so that only $h_+$ contributes, the other polarization not. It was derived at the Jet Propulsion Laboratory, see below. The coefficients, as the limits below show, are such that: for $\sin\theta=1$ the first coefficient vanishes, the other two are $2h_+$-terms with the same argument and cancel; for $\theta=0$ the third coefficient vanishes.)

\begin{coursenote}[Which angle? the geometry of the limits]
The lecturer has no diagram at first and draws it later: a laser sending light, a mirror at distance $L$, a plane wave with wave fronts; ``$\theta$ is the angle between the wave fronts and the path of the light''; and then says that ``$\theta=\pi/2$ is the wave parallel to the light, going along the same direction as the light''. Read together these two sentences are consistent only if $\theta$ is measured so that $\theta=0$ is the broadside incidence (wave fronts along the light path: wave perpendicular to the light) and $\theta=\pi/2$ the wave travelling along the light. The limits below follow that convention \ts{28}{1}{16:31}\ts{28}{1}{16:51}\ts{28}{1}{17:12}\ts{28}{1}{17:33}.
\end{coursenote}

\paragraph{Limits.}
\begin{itemize}
  \item $\theta=\pi/2$ (wave and light parallel, $\sin\theta=1$): the first term has a factor zero, the second a factor two $h_+(t)$, the third a factor two with the argument $t+L(1-1)=t$, the same as in the other: the two cancel and the result is zero: ``when a gravitational wave goes along the light you don't get any disturbance on the light at all, because they keep phase with each other: the action is all transverse'' \ts{28}{1}{17:54}\ts{28}{1}{18:14}\ts{28}{1}{18:34}.
  \item $\theta=0$ (broadside, ``where you expect the maximum effect''): the first and the second term fully present, the third absent: only the amplitude of the wave at emission and at return matters, \emph{not even} the time of bounce off the mirror \ts{28}{1}{18:54}\ts{28}{1}{19:15}.
  \item In between the mirror term has some effect \ts{28}{1}{19:35}.
\end{itemize}
It is a measurable derived from general relativity in any coordinate system, assuming only that the gravitational wave travels at the speed of light; it has nothing to do with the \emph{wavelength of the light}. ``People ask: you do an interferometry experiment, is the wavelength of light stretched? That is irrelevant: it is all to do with the return time of the light going back and forth'' \ts{28}{1}{19:35}\ts{28}{1}{19:56}\ts{28}{1}{20:16}.

\section{Kinds of beam detectors}\label{sec:c28-kinds}

\paragraph{Spacecraft tracking (Doppler tracking).} Since the 1960s people at NASA understood that sending signals out to spacecraft in the solar system and monitoring the return time performs exactly the three term experiment, and one might detect gravitational waves that way. The three term formula was derived by people at the Jet Propulsion Laboratory in California for NASA: the first time the quantitative interaction of gravitational waves with light was fully explored. First done in the late 1960s and early 1970s (the Pioneer missions among them), many later missions to the outer planets carried such experiments: radio waves are sent out and \emph{transponded} (a transponder is a radio receiver--transmitter which answers a signal with minimum or known delay) \ts{28}{1}{20:36}\ts{28}{1}{20:57}\ts{28}{1}{21:18}\ts{28}{1}{21:38}.

The sensitivity is limited because one is not in vacuum: the \emph{interplanetary plasma} affects the travel time of radio waves, so the return time is not only a function of gravitational waves (noise from the plasma); and one needs a very good clock to measure the return time: the atomic clocks of the Deep Space Network, some of the best in the world, stable to about $10^{-16}$ (the lecturer: ``$10^{-16}$ or $10^{-17}$, well probably $10^{-16}$''); that is the limit. Typical $h\sim10^{-16}$, in the millihertz band (set by the time scale of the transponding). Many runs, the most recent the Galileo mission: ``numbers like that, but they have never detected gravitational waves'' \ts{28}{1}{21:59}\ts{28}{1}{22:20}\ts{28}{1}{22:40}\ts{28}{1}{23:01}\ts{28}{1}{23:22}.

\paragraph{Pulsar timing.} A variation: pulses of radio waves from distant pulsars travel through the gravitational wave field and arrive at the Earth. Nothing is sent back to the pulsar (thousands of light years away), so it is one way, not up and back; but the one-way calculation gives similar formulas. $h$ matters both at the Earth and at the pulsar: the field passing the pulsar a thousand years ago when the pulse was emitted, and the field passing the Earth now. The pulsars are all over the sky and the fields at the different pulsars are not correlated with one another, which is basically noise: averaging over all pulsars one keeps only the correlated effect, the gravitational wave at the Earth at the time the pulses arrive \ts{28}{1}{23:42}\ts{28}{1}{24:03}\ts{28}{1}{24:24}\ts{28}{1}{24:44}\ts{28}{1}{25:05}.

The sensitivity is limited by clocks on Earth and by intrinsic pulsar irregularities. \textbf{Millisecond pulsars} are extremely clean: the individual pulses vary, but averaged over about a year, after taking out a few systematic trends, one gets a very clean signal; about $20$ pulsars around the sky could be used. They are so clean that even with the best atomic clocks the irregularities cannot be detected: on time scales of about a year pulsars may be better clocks than the best clocks that can be built. So the limits are set by the clocks, and the pulsar irregularities are on the same scale or better; as clocks improve pulsar timing improves; at present $\sim10^{-16}$. The predictions of Sesana for the backgrounds (previous lecture) are in this range, around $10^{-16}$, ``strong, chaotic, stochastic signals'': so, unlike spacecraft tracking, pulsar timing has a chance of detecting even with today's clocks. The frequency band is not set by the transmission time (thousands of years) but by the averaging time, about a year: \textbf{one to ten nanohertz}, below the LISA band, ``a totally different frequency band than anything else accessible to us'', with supermassive black hole binaries as sources \ts{28}{1}{25:25}\ts{28}{1}{25:47}\ts{28}{1}{26:08}\ts{28}{1}{26:29}\ts{28}{1}{26:49}\ts{28}{1}{27:09}\ts{28}{1}{27:30}.

\paragraph{A one arm detector with a clock.} Why not build a detector with one arm and a very good clock, on the ground or in space, like the thought experiment? Interferometers with two arms need big vacuum systems and civil engineering, expensive. The one arm would be ``the instrument of choice if you could do it'', but it is \emph{clock limited}, now $10^{-16}$. The people who build atomic clocks hope for $10^{-18}$ in the very near future, $10^{-19}$ or $10^{-20}$ maybe; but even $10^{-20}$ is not good enough for the sensitivity needed for ground-based detection in the kilohertz band, so interferometers are needed. In principle the frequency band would be determined by the arm length (the time up and down is a sampling time) and by the clock stability time scale \ts{28}{1}{27:50}\ts{28}{1}{28:11}\ts{28}{1}{28:32}\ts{28}{1}{28:52}\ts{28}{1}{29:13}.

\begin{coursenote}[Student question: clocks in space?]
A student asks whether this means putting very sensitive clocks on spacecraft that go far away (and whether there is anything like it on Voyager or Pioneer). The lecturer: on the ground one would just have a clock where the laser is sent out; in space one has the additional complication of a really good atomic clock; and the clocks in space are ``always one or two generations behind the ones you can build on the ground'' \ts{28}{1}{29:34}\ts{28}{1}{29:54}.
\end{coursenote}

\paragraph{Atom interferometry.} Particles are waves at the quantum level (wave--particle duality) and an interferometer can interfere atoms, not light, as has been done to demonstrate the wave nature of matter; one could in principle detect gravitational waves this way (``you would have a different three term relation''), though in the far future and with no practical way known. The people of atom interferometry propose a space one arm detector in which the atom interferometer takes the role of the clocks: clouds of atoms are ejected at both ends, one passes a beam of light through the atoms at one end imprinting a quantum state (``it is all to do with entangled states''), at the other end one resolves the entanglement: hopefully a higher accuracy than atomic clocks; one measures a phase or transmission difference, not time with a full clock. The lecturer sits on an ESA committee discussing it: it does not look like a technology for the time scale of LISA but might come for a next generation; ``it would be very sexy'' \ts{28}{1}{30:15}\ts{28}{1}{30:36}\ts{28}{1}{30:56}\ts{28}{1}{31:17}\ts{28}{1}{31:38}\ts{28}{1}{31:58}\ts{28}{1}{32:19}\ts{28}{1}{32:39}.

\section{The laser interferometer}\label{sec:c28-interferometer}

Why the laser interferometer with two arms and not one? Physically the second arm is the \textbf{time reference} for the first: one cannot build a clock accurate to that level, but with two cavities both being stabilised one beats one cavity against the other; both cavities are affected by the gravitational wave as it comes through, so one picks up the \emph{difference} between what happens in the two arms. The three term relation has to be put together in two perpendicular ways (with different angles $\theta$), but the response is fairly straightforward to work out \ts{28}{1}{32:59}\ts{28}{1}{33:20}\ts{28}{1}{33:41}\ts{28}{1}{34:01}\ts{28}{1}{34:23}.

The sensitivity is limited by how well the cavities can be built: environmental noise, seismic noise, laser noise, photon shot noise, thermal effects in the mirrors, the quality and shape of the mirrors. One believes one can get $10^{-22}$ in the next few months, $10^{-23}$ or $10^{-24}$ when the technologies are pushed further \ts{28}{1}{34:44}\ts{28}{1}{35:04}.

\paragraph{The frequency band.} Determined basically by the arm length but not equal to the arm length or its inverse: multiple reflections build up the signal in a cavity which stays much longer than a single light crossing time; the cavities may have a natural frequency of about $100$ kHz but one measures gravitational waves of $1$ kHz or $100$ Hz by keeping the light in the cavities for much longer. In space, with arms of millions of kilometres instead of a few, one goes to the \emph{millihertz} band \ts{28}{1}{35:24}\ts{28}{1}{35:45}.

\subsection{The short arm limit}

No assumption has been made on the length $L$ in relation to the wavelength of the gravitational waves. If $L$ is short compared to the wavelength, $2L$ is a short time compared to the period, and one can Taylor expand $h$ and keep only first order in $L$ \ts{28}{1}{36:06}\ts{28}{1}{36:27}\ts{28}{1}{36:47}. Setting $L=0$ all terms cancel (a no arm detector gives no signal), so to first order the signal is \emph{proportional to $L$}. Physically, for a detector short compared to the wavelength the tidal effect is coherent across the whole detector, which is just stretched and compressed in proper distance, depending only on $L$: the simple picture, valid on the ground \ts{28}{1}{37:07}\ts{28}{1}{37:28}\ts{28}{1}{37:48}\ts{28}{1}{38:09}.

The picture fails if $L$ is comparable to or bigger than the wavelength: the light goes through a fluctuating wave field and there are cancellations; ``all of that is in this formula, you can rely on it'', but it is not a simple coherent expansion and contraction. LISA is in that regime: arms of millions of kilometres are longer than or comparable to some of the wavelengths, which affects its sensitivity at the high frequency end \ts{28}{1}{38:09}\ts{28}{1}{38:29}\ts{28}{1}{38:50}.

\section{The cosmic microwave background}\label{sec:c28-cmb}

Another way, quite different: the BICEP2 excitement of the previous year (2014) about the signal of gravitational waves. It is \emph{not} a beam detector: one does not see time variations in the transmission of the microwaves, but the \emph{imprint of gravitational waves on the polarization} \ts{28}{1}{39:11}\ts{28}{1}{39:31}\ts{28}{1}{39:52}. A photon scattering off charges which move has its polarization altered: if the charge is stationary the polarization is random, but if the plasma moves in one direction there is preferential (transverse) polarization. A gravitational wave coherently stretching the plasma just when the universe becomes transparent leaves an imprint in the polarization of the microwaves: ``a bit more like a bar experiment'': a mechanical system distorted by the wave, and the microwaves carry the signal of what happened \ts{28}{1}{40:13}\ts{28}{1}{40:34}\ts{28}{1}{40:55}.

The density perturbations were about one part in $10^6$ (the nice fluctuation diagrams of the CMB), so the gravitational wave perturbations might have a similar amplitude; BICEP2 would allow detecting that amplitude if present. The angular resolution of the CMB experiments (each cell of the sky) determines the wavelengths: one sees in each cell a coherent gravitational wave signal whose wavelength at the time fills that angular region, about a thousandth of the Hubble horizon at that time \ts{28}{1}{41:16}\ts{28}{1}{41:37}\ts{28}{1}{41:59}.

As the team had to agree, measurements from the Planck satellite indicate that BICEP2 did not have the sensitivity to see gravitational waves: they reflect foreground fluctuations (caused by intermediate ionization and the plasma) which BICEP2 could not exclude. One may peer through since the foregrounds have a different \emph{frequency dependence} from the microwave background: BICEP2 measured one frequency; measuring over several one can subtract the foreground and get below it, which people are doing now. If the gravitational waves are at the $10^{-6}$ level one may have a measurement in about a year; but many theorists predict the signal is actually zero, so upper limits are possible, and it is hard to predict when that level is reached \ts{28}{1}{42:21}\ts{28}{1}{42:42}\ts{28}{1}{43:03}\ts{28}{1}{43:23}.

\section{Building a laser interferometer}\label{sec:c28-building}

Interferometers are wonderful instruments because they involve practically every branch of physics: thermodynamics (controlling the thermal effects in the mirrors), materials (suspending and constructing the mirrors), optics, quantum optics \ts{28}{1}{43:44}\ts{28}{1}{44:06}.

\paragraph{The two main limitations.} Environmental disturbances and laser shot noise \ts{28}{1}{44:29}.
\begin{itemize}
  \item \textbf{Environmental noise (ground)}: removed by \textbf{pendulum suspensions}: the mirrors hang from a metre or more of suspension. A pendulum is a mechanical \emph{low pass filter}: moving the top at a frequency high compared with the pendulum frequency, the pendulum hardly moves; at low frequency it follows. LIGO uses a three stage suspension (three ganged pendulums, developed in GEO), the product of three filter functions, a very steep falloff of seismic noise; the turning point is around a hertz, so there is a lot of noise at a hertz and it falls off fast: by around $10$ Hz one hopes it is low enough. The \emph{lower limit} of the frequency band is set by the pendulum. Virgo does an even better job with six stage suspensions \ts{28}{1}{44:49}\ts{28}{1}{45:10}\ts{28}{1}{45:31}\ts{28}{1}{45:51}\ts{28}{1}{47:15}.
  \item \textbf{Laser shot noise}: at the upper end few noise sources come in; the limit is the amount of laser light, since light is quantized: the photons arrive at the beam splitter in a random way and give tiny fluctuations of the intensity of the light coming out, even in a perfectly stable interference pattern. This dominates at high frequencies. ``If you understand these two kinds of noise you understand the sensitivity of gravitational wave detectors pretty much.'' One attacks the shot noise by \textbf{squeezing} (below), and the radiation pressure trade-off (see the questions) \ts{28}{1}{46:12}\ts{28}{1}{46:33}\ts{28}{1}{46:54}\ts{28}{1}{47:35}\ts{28}{1}{47:56}.
\end{itemize}

\paragraph{The Michelson interferometer.} Laser, beam splitter, two mirrors (one up, one over), a photodetector. The light is split into two beams which reflect and come back; at the beam splitter part of the light goes back towards the laser (constructive pattern) and part goes out to the detector (destructive pattern). Normally one runs so that nothing comes out there: when there is a signal one goes from zero to something, a good signal to noise, limited only by the noise of the photodetector. The disturbing consequence: all the light goes back towards the laser, hits the wall and ``burns a hole in your wall'': lost \ts{28}{1}{48:16}\ts{28}{1}{48:37}\ts{28}{1}{48:57}.

\paragraph{Storing the light.} Why bounce only once up and down? One can store the light in the arms for longer than the light travel time, which effectively increases the arm length: in GEO there is one bounce (a folded arm); in a prototype built at Garching near Munich the light made about $40$ independent bounces on spots of a big mirror, elaborate and difficult to control, ``we don't do it that way any more''; the principle is the same. Since one is in the short arm regime, the wave is coherent over the whole detector, and if one stores the light for less than about a quarter of the period of the gravitational wave one gets a coherent effect on all the bounces, and the signal builds up: ``it effectively makes the arm length much longer than you could build in practice'' \ts{28}{1}{49:18}\ts{28}{1}{49:43}\ts{28}{1}{50:03}\ts{28}{1}{50:24}.

\paragraph{Power recycling.} Ron Drever at Glasgow had the idea: catch the light coming back, the wasted light, and send it back in, in phase with the new light (the gravitational wave signal is not in that light: it is in the difference). Proposed independently, at about the same time, by Schilling in Munich [name as heard]. With perfect mirrors one would get ``a laser bomb'': the light builds up until it is limited by the losses at the mirrors (reflection, scattering, absorption), and the laser only supplies the losses. With mirror losses of one part in $10^{6}$, one has about \emph{a million times} as much power inside the cavities as the laser power coming in: a huge improvement in sensitivity, a huge reduction of the shot noise \ts{28}{1}{50:44}\ts{28}{1}{51:04}\ts{28}{1}{51:25}\ts{28}{1}{51:47}\ts{28}{1}{52:07}\ts{28}{1}{52:27}\ts{28}{1}{52:48}.

\begin{coursenote}[Student question: why not remove the laser's own mirrors?]
A student suggests removing the mirror inside the laser box: ``no, because then a huge amount of light would go through the laser and distort all the optics; you want a really sharp line there and the laser very well controlled; it is much better to adjust the phase with feedback systems so that they are in phase, which is pretty easy''. Remark: GEO has only one bounce but gets all its sensitivity through the light recycling; the power recycling light then goes through the beam splitter with a huge thermal load, a limitation on GEO's sensitivity; LIGO and Virgo, with many more reflections, store the power in the arms and put less load on the beam splitter, an advantage: ``technical things'' \ts{28}{1}{53:08}\ts{28}{1}{53:30}\ts{28}{1}{53:50}\ts{28}{1}{54:12}\ts{28}{1}{54:32}.
\end{coursenote}

\paragraph{Signal recycling.} The idea first proposed in the late 1980s in the Glasgow group: the lecturer cannot remember the name of the inventor (``unfortunately killed in a climbing accident'') and corrects his first attribution to Drever. A mirror is put at the output port where the fringes are: with no gravitational wave it does nothing; with a wave one can tune it, make it partly reflecting and place it so as to tune the whole cavity and reflect the signal back and forth, building up the signal in a particular bandwidth: the higher the reflectivity the narrower the band, ``a kind of resonant detector with narrow bandwidth''; with partial reflectivity the signal is shaped. It has been done and tested in GEO and will be done in the big detectors, mainly to shape the signal, improving the sensitivity in particular frequency bands for a particular signal \ts{28}{1}{54:52}\ts{28}{1}{55:13}\ts{28}{1}{55:34}\ts{28}{1}{55:54}\ts{28}{1}{56:15}\ts{28}{1}{56:35}.

\paragraph{Arm cavities.} In LIGO and Virgo, a whole \emph{cavity} in each arm instead of a bunch of multiple reflections: the input mirror is partially transmitting (the lecturer: ``$99$ percent silver to something'') and the end one totally reflecting, and most of the power in the arms is stored in these cavities: this is how most of these detectors work now \ts{28}{1}{56:58}\ts{28}{1}{57:20}.

\section{Two misconceptions}\label{sec:c28-myths}

``Two things I have been saying to you all along, which you probably believe, are wrong, and now you can understand why'' \ts{28}{1}{57:20}\ts{28}{1}{57:42}.

\begin{enumerate}
  \item \emph{``We need freely falling test masses, which is why the mirrors are hung.''} It is not the reason: the mirrors are hung to isolate them from ground vibration, by mechanical filtering. They would be free anyway, even nailed down on a rigid suspension attached to the ground: the Earth is not rigid enough to stop the motion. The gravitational wave acts on a time scale of milliseconds; the rigidity of the Earth depends on pressure waves going through it to counteract the stretching, and seismic waves are rather slow: when one station moves, another a couple of kilometres away does not feel the pull until much later than a few milliseconds. So for gravitational waves the two stations are free to move anyway, in the horizontal direction and in any direction \ts{28}{1}{57:42}\ts{28}{1}{58:03}\ts{28}{1}{58:23}\ts{28}{1}{58:45}\ts{28}{1}{59:06}. The detectors are sensitive in the horizontal direction only, because the suspension isolation carries them up and down: ``we don't detect gravitational waves in the vertical direction; it is all to do with the suspension and vibration isolation, not with the need for freely falling bodies; they are freely falling bodies on this time scale'' \ts{28}{1}{59:26}.
  \item \emph{``Bigger is better: we build kilometre sized detectors to build up the signal.''} Wrong, because of the multiple reflections which increase the effective arm length: one could build GEO600 at $600$ m with the same signal sensitivity as a $3$ km one, with more reflections and better optics. Why bigger is better: other noise sources, particularly the vibration noise coming through the suspension and the \emph{thermal vibrations in the mirrors}: the more times the light bumps into a mirror the more noise is picked up; so one wants the mirrors far apart to store light in the arms while minimising the number of contacts with the mirrors (GEO's light touches the mirrors four or five times more often than in the larger detectors: a background of noise much harder to get rid of). With a $10$ km underground detector like the Einstein Telescope the effect is reduced even more: ``the reason for going to long arms is to minimise the mechanical noise inputs when you are contacting the mirrors'' \ts{28}{1}{59:47}\ts{28}{1}{1:00:07}\ts{28}{1}{1:00:28}\ts{28}{1}{1:00:48}\ts{28}{1}{1:01:08}.
\end{enumerate}

\section{Quantum noise and squeezed light}\label{sec:c28-squeezing}

The laser interferometer diagram, here as GEO (one reflection, a folded arm, all the light through the beam splitter and the power on the power recycling mirror; a signal recycling mirror not shown). It is a closed system except at one place: the output port, where the signal comes out, but also where \emph{the vacuum comes in}: the electromagnetic vacuum, which at a quantum level is noisy, with the zero point fluctuations of all the quantum modes of the field. The shot noise (photon arrival times) can be formulated entirely as these vacuum fluctuations entering the detector, imposed as noise on a smooth classical field \ts{28}{1}{1:01:30}\ts{28}{1}{1:01:52}\ts{28}{1}{1:02:12}\ts{28}{1}{1:02:34}\ts{28}{1}{1:02:54}.

To control it: put an apparatus and a mirror at that port, stop the vacuum coming in ``in a raw way'' and \textbf{squeeze} it: turn it into a different vacuum state with correlations. ``You are not getting around the Heisenberg uncertainty principle; you are changing the conjugate variables where the uncertainty applies and squeezing the uncertainty in one of them: you put all the noise from the vacuum into something you do not measure.'' One measures the variable (quadrature) whose noise is squeezed, and makes sure that it contains the gravitational wave signal: ``we do not want to make a quantum measurement, we want to get rid of the quantum noise into one quadrature and have in the other a very quantum quiet background'' \ts{28}{1}{1:03:14}\ts{28}{1}{1:03:36}\ts{28}{1}{1:03:56}\ts{28}{1}{1:04:17}\ts{28}{1}{1:04:37}.

In GEO600 it works: on the sensitivity curve, the black line before squeezing, after squeezing the noise is reduced at the high frequencies where shot noise dominates (not at the low ones, dominated by vibration noise), by a factor of two, at most two and a half or three. A factor two in noise improves the reach by two and the surveyed volume by eight: ``you get ten times as many gravitational wave sources by just introducing squeezing''. It is all done with ``the deepest quantum mechanics, totally mysterious''; the apparatus (the table with the squeezed light apparatus next to the GEO vacuum tanks) is not trivial. It is being put into the big interferometers in a second stage; one will get below the standard shot noise limit, but ``it will still be laser noise, photon noise'' \ts{28}{1}{1:04:58}\ts{28}{1}{1:05:19}\ts{28}{1}{1:05:40}\ts{28}{1}{1:06:00}\ts{28}{1}{1:06:21}.

\section{LISA Pathfinder and the transponder}\label{sec:c28-lisa}

Pictures of LISA Pathfinder (the lecturer first mis-clicks the slides): a big can with two test masses half a metre apart, the free test masses, the interferometry along the optical bench, ``very similar to what we will do with full LISA, just the separation isn't there'': not a gravitational wave detector but a good test of the principles of the interferometry. LISA has no cavities: one sends a signal up and it comes back, measuring the return time. But a light beam focused as narrowly as possible still spreads over a million kilometres to tens of kilometres wide, compared to the receiving spacecraft, so there is a big reduction in the received light; with one watt lasers the received light is very small, and if just reflected back by a mirror ``we would never get any photons back at all'' \ts{28}{1}{1:06:42}\ts{28}{1}{1:07:03}\ts{28}{1}{1:07:24}\ts{28}{1}{1:07:45}\ts{28}{1}{1:08:05}\ts{28}{1}{1:08:27}.

So, as the spacecraft people do, a \textbf{transponder}: a receiving system senses the incoming light and an independent laser, slaved to the incoming light, sends back with the right phase: ``an active, amplifying mirror, the heart of the LISA technology''. In the triangle, any two arms can be put together to form an interference pattern, giving three independent gravitational wave signals, really independent if the waves are short compared to the arms because they affect the arms differently: more information about the waves. ``Basically the technology is that of Pathfinder plus the transponding to get the right intensity'' \ts{28}{1}{1:08:47}\ts{28}{1}{1:09:07}\ts{28}{1}{1:09:27}\ts{28}{1}{1:09:48}.

\section{Questions}\label{sec:c28-questions}

\begin{coursenote}[Student questions]
(1) \emph{The wavelength.} How can the wavelength not matter for the return time when one sees an interferometer? Answer: one can do white light interferometry ``if I am really good'': with exactly equal arms every wavelength comes back and interferes with itself in the right way, so the wavelength is not relevant. The interferometer measures \emph{differences in return time}: if the transmission time along one arm changes relative to the other the signal is out of phase. The wavelength matters for how accurately one can measure it, but not for the effect of the gravitational wave imprinted on it \ts{28}{1}{1:10:30}\ts{28}{1}{1:10:50}\ts{28}{1}{1:11:11}\ts{28}{1}{1:11:32}.

(2) \emph{Radiation pressure.} With many photons between the mirrors, wouldn't the pressure of the photons already move the mirrors? Yes, there is a balance: photon pressure from many photons is traded off against the accuracy with which one can measure the shot noise. At present one is limited by shot noise, but when the laser power is pumped to the final levels (about $100$ or $200$ W input) one will be in a trade-off situation; there are many more optical tricks, squeezing among them \ts{28}{1}{1:11:52}\ts{28}{1}{1:12:13}\ts{28}{1}{1:12:33}.

(3) \emph{Signal recycling and shot noise.} Does signal recycling increase the shot noise since less light goes to the detector? No: because of the cavity roughly the same amount of light comes out, and the shot noise is in the interference pattern at the beam splitter, ``it is the randomness of the interference, not the photon counting at the photodiode'' \ts{28}{1}{1:12:55}\ts{28}{1}{1:13:15}.

(4) \emph{Squeezing once more.} In the standard description of the electromagnetic field one builds a Hamiltonian on the number of photons, with an uncertainty between the number of photons and the time of arrival (a student corrects: it is the \emph{phase}, ``yes, that's right, thank you''). But one need not take these as the conjugate variables: one can go over to a description in which the two conjugate variables are the two outputs from the beam splitter, the intensities of the constructive and destructive interference, which depend on the phase shift between the two arms. Then the question is whether one can find the technology to change the Heisenberg uncertainty circle into an ellipse so that the uncertainty in the intensity of light going out to the detection point is reduced, at the cost of increasing the uncertainty of the light going back to the power recycling mirror: ``you try to choose variables that correspond to the two different exit paths of the interferometer; how you do that I am at a loss to describe myself''; the principle is to make the part of the light carrying the signal as quiet as possible, ``I don't get around the uncertainty, I just put more uncertainty in one conjugate variable and less in the other'' \ts{28}{1}{1:13:36}\ts{28}{1}{1:13:57}\ts{28}{1}{1:14:26}\ts{28}{1}{1:14:47}\ts{28}{1}{1:15:08}\ts{28}{1}{1:15:30}\ts{28}{1}{1:15:50}\ts{28}{1}{1:16:11}\ts{28}{1}{1:16:32}\ts{28}{1}{1:16:52}\ts{28}{1}{1:17:12}.

(5) \emph{A really long optical fibre inside a satellite?} (``perhaps a stupid question'': it is not.) Light in a fibre does not go at the speed of light; fibres are intrinsically imperfect, with a lot of scattering compared to a vacuum cavity, so the light one can store is limited; and there are thermal effects in the materials: even with no external disturbance one gets little poppings, as seen with the coatings of the mirrors, where turning up the light intensity gives many glitches thought to come from thermal effects (little domains out of equilibrium releasing or changing); fibres as drawn do this a lot. So at the moment a vacuum cavity is much better than a material system containing the light \ts{28}{1}{1:17:36}\ts{28}{1}{1:17:57}\ts{28}{1}{1:18:18}\ts{28}{1}{1:18:38}\ts{28}{1}{1:18:59}.

(6) \emph{A mirror on the Moon as an obviously very long one arm instrument?} Yes, the people who do laser ranging to the Moon have talked of setting limits on gravitational waves in that frequency band, but the reflecting area of those mirrors is very small, so one counts individual photons: photon shot noise dominates and limits the sensitivity a lot; huge mirrors on the Moon would help, but there is also the laser light travelling through the atmosphere and other noise imprinted on it \ts{28}{1}{1:19:25}\ts{28}{1}{1:19:47}\ts{28}{1}{1:20:08}\ts{28}{1}{1:20:29}.
\end{coursenote}

\begin{aside}[The end of the school]
The lecture ends the school. The organiser thanks the lecturer for ``such a masterful conclusion to the lectures'', says that the students complimented his lectures already in the break and will remember the school for a long time, and gives him a cup to remember the lectures by and, ``a happy coincidence that your German is so good'', a copy of Kepler's \emph{Astronomia nova} with many comments, ``very highly recommended''; the lecturer thanks him \ts{28}{1}{1:20:57}\ts{28}{1}{1:21:18}\ts{28}{1}{1:21:38}.
\end{aside}

\section*{Key concepts}
\begin{keyconcepts}
\textbf{Retrieval questions.}
\begin{enumerate}
  \item How does a bar detector work, why is it limited, and what is its relation to the resonant frequency? What happened in the Weber controversy and why does it still matter?
  \item Why is the return time of a light signal in free fall a gauge independent measurable? Why does the wavelength of the light not matter?
  \item State the structure of the three term relation: which three times of the waveform enter? What happens for a wave travelling along the light, and for broadside incidence?
  \item Compare spacecraft tracking, pulsar timing, and a one arm detector with a clock: what limits each, and which frequency band does each have?
  \item Why does a two arm interferometer work where a one arm one needs a clock? What is the short arm limit and where does it fail?
  \item How does BICEP2 look for gravitational waves, and why is it not a beam detector?
  \item What are the two main noise sources of an interferometer, and which tricks address them (pendulum suspensions, stored light, power recycling, signal recycling, squeezing)?
  \item Why are the mirrors hung, and why are the detectors kilometres long? (the two misconceptions)
  \item What does squeezed light do to the Heisenberg uncertainty, and how much does a factor two in noise gain in volume?
  \item What is a transponder and why does LISA need one?
\end{enumerate}
\textbf{Concepts.}
\begin{itemize}
  \item Bar detectors and the Weber controversy (\cref{sec:c28-bars}).
  \item Beam detectors, return time, the three term relation (\cref{sec:c28-beam}, \cref{thm:c28-three}).
  \item Spacecraft tracking, pulsar timing, one arm with clocks, atom interferometry (\cref{sec:c28-kinds}).
  \item Laser interferometer, short arm limit (\cref{sec:c28-interferometer}).
  \item The cosmic microwave background polarization (\cref{sec:c28-cmb}).
  \item Noise, suspensions, light storage, recycling, the two misconceptions (\cref{sec:c28-building}, \cref{sec:c28-myths}).
  \item Squeezed light (\cref{sec:c28-squeezing}); LISA Pathfinder and the transponder (\cref{sec:c28-lisa}).
\end{itemize}
\end{keyconcepts}

\flushside
\end{paracol}
